\section*{Chapitre \ref{ch:b1:rifts-ridges-margins} --- Rifts, dorsales et marges passives}

\begin{solution}{exo:b1:rifts-ridges-margins:1}
(a) Le dyke couvrirait $\qty{60000}{m}\times\qty{8}{m} = \qty{4.8e5}{m^2}$ en
plan, donc sa hauteur moyenne serait $\qty{2.5e9}{m^3}/\qty{4.8e5}{m^2} =
\qty{5200}{m}$, soit environ \qty{5}{km}~: le dyke traverserait une grande part de la \omterm{def:b1:layered-earth:layers}{croûte}.
(b) $\qty{8}{m}/\qty{0.018}{m.\year^{-1}} = 444$ ans, environ quatre siècles.
\end{solution}

\begin{solution}{exo:b1:rifts-ridges-margins:2}
(a) Le \omterm{def:b1:cross-sections:throw}{rejet horizontal} est le \omterm{def:b1:cross-sections:throw}{rejet vertical} divisé par la tangente du \omterm{def:b1:geological-maps:attitude}{pendage}~:
$3.0/\tan\ang{60} = \qty{1.73}{km}$ et $0.8/\tan\ang{70} = \qty{0.29}{km}$~;
l'extension de la coupe est de \qty{2.02}{km}. (b) Avant la fracturation la
coupe était large de $25 - 2.02 = \qty{22.98}{km}$, soit environ \qty{23}{km}~: elle s'est
allongée d'environ \qty{9}{\percent}.
\end{solution}

\begin{solution}{exo:b1:rifts-ridges-margins:3}
Ultralentes (moins de \qty{20}{mm.an^{-1}})~: la dorsale de Gakkel (12), le centre de la mer
Rouge (15) et la dorsale de Reykjanes (19), cette dernière proche de la limite. Lentes (20 à
\qty{55}{mm.an^{-1}})~: la dorsale de Carlsberg occidentale (23) et l'Atlantique sud
(33). Rapide (au-dessus de \qty{75}{mm.an^{-1}})~: la dorsale est-Pacifique (143). Les
dorsales ultralentes et lentes devraient avoir une vallée axiale entre des failles normales~; la
dorsale est-Pacifique un haut axial.
\end{solution}

\begin{solution}{exo:b1:rifts-ridges-margins:4}
(a) $\qty{3.0}{km^2}\times\qty{7.1}{km} = \qty{21.3}{km^3}$, soit environ
\qty{21}{km^3} par an. (b) $\qty{3.028e8}{km^2}/\qty{3.0}{km^2.\year^{-1}} =
\qty{1.01e8}{\year}$, \qty{101}{\Ma}. Le plus vieux plancher, environ \qty{190}{\Ma}, est
plus vieux parce que le plancher n'est pas détruit par ordre d'âge~: la subduction consomme
du plancher de tous âges, du vieux plancher survit là où aucune fosse ne l'a atteint,
et le taux de production a varié.
\end{solution}

\begin{solution}{exo:b1:rifts-ridges-margins:5}
(a) $2500 + 350\times2 = \qty{3200}{m}$, $2500 + 350\times6 = \qty{4600}{m}$,
$2500 + 350\times8 = \qty{5300}{m}$. (b) $t = \bigl((d - 2500)/350\bigr)^2$~:
$(2100/350)^2 = \qty{36}{\Ma}$ et $(2800/350)^2 = \qty{64}{\Ma}$. (c) Le
plancher s'enfonce de \qty{1400}{m} entre 4 et \qty{36}{\Ma}, mais de \qty{700}{m} seulement
entre 36 et \qty{64}{\Ma}, sur une durée comparable.
\end{solution}

\begin{solution}{exo:b1:rifts-ridges-margins:6}
Le coefficient de dilatation effectif est
$\num{3.2e-5}\times3300/2300 = \qty{4.591e-5}{K^{-1}}$, et
$\sqrt{\kappa\times\qty{1}{\Ma}} = \sqrt{\num{e-6}\times\num{3.156e13}} =
\qty{5618}{m}$. Alors $C = (2/\sqrt\pi)\times\num{4.591e-5}\times1350\times5618 =
1.1284\times\num{4.591e-5}\times1350\times5618 = \qty{393}{m.\Ma^{-1/2}}$. Pour un
\omterm{def:b1:layered-earth:layers}{manteau} plus froid de \qty{250}{K}, $C$ varie comme $T_m - T_0$~:
$393\times1100/1350 = \qty{320}{m.\Ma^{-1/2}}$. Le plancher d'un océan plus froid
s'enfonce plus lentement et suit la plus petite constante. La valeur ajustée de \num{350}
du chapitre n'en est pas un exemple~: elle reste dans l'incertitude des paramètres.
\end{solution}

\begin{solution}{exo:b1:rifts-ridges-margins:7}
(a) $\qty{e-6}{cal.cm^{-2}.s^{-1}} = \num{4.184e-6}\times\num{e4} =
\qty{4.184e-2}{W.m^{-2}} = \qty{41.84}{mW.m^{-2}}$, donc la constante est
$11.3\times41.84 = \qty{473}{mW.m^{-2}.\Ma^{1/2}}$, \qty{3.5}{\percent} sous
les 490 du \cref{ch:b1:heat-budget}~: les deux ajustements concordent. (b) $q = 473$,
$473/\sqrt{10} = 150$ et $473/\sqrt{150} = \qty{38.6}{mW.m^{-2}}$. La loi en
racine carrée continue de baisser, alors que le modèle de plaque se stabilise à une constante~:
à \qty{150}{\Ma} la valeur de la racine carrée est trop basse, comme la loi de profondeur va trop
profond.
\end{solution}

\begin{solution}{exo:b1:rifts-ridges-margins:8}
Prenons le segment nord sur $x = 0$ et le segment sud sur $x = -40$
(km), la transformante sur $y = 0$ entre eux. (a) Juste au nord de la transformante le
plancher est à l'ouest du segment nord et se déplace vers l'ouest~; juste au sud, il
est à l'est du segment sud et se déplace vers l'est. Vu de l'un ou l'autre côté, le
côté opposé se déplace vers la gauche~: le glissement est senestre. (b) La
dorsale semble avoir été déplacée de façon dextre, sa partie sud vers
l'ouest~: le sens opposé, comme le prédit la \cref{prop:b1:rifts-ridges-margins:transform}.
(c) Les séismes ne se produisent que sur les \qty{40}{km} entre les deux extrémités de dorsale,
là où les deux côtés appartiennent à des plaques différentes~; au-delà, le long des
\omterm{def:b1:rifts-ridges-margins:transform}{zones de fracture}, les deux côtés appartiennent à la même plaque et se déplacent ensemble.
\end{solution}

\begin{solution}{exo:b1:rifts-ridges-margins:9}
La demi-vitesse est de \qty{10.6}{mm.an^{-1}}, \qty{10.6}{km.\Ma^{-1}}. (a) Le saut
d'âge est $350/10.6 = \qty{33.0}{\Ma}$. (b) À l'extrémité de la transformante, un plancher
de \qty{33.0}{\Ma} ($2500 + 350\sqrt{33.0} = \qty{4511}{m}$) fait face à un plancher
d'âge 0 (\qty{2500}{m})~: un ressaut d'environ \qty{2010}{m}. \qty{100}{km} au-delà,
les âges sont $450/10.6 = 42.5$ et $100/10.6 = \qty{9.4}{\Ma}$, les profondeurs
4781 et \qty{3575}{m}~: un ressaut d'environ \qty{1210}{m}. \qty{1000}{km} au-delà,
les âges sont 127.4 et \qty{94.3}{\Ma}, les profondeurs 6450 et \qty{5899}{m}~:
\qty{550}{m}. (Ces âges dépassent \qty{80}{\Ma}, où la loi en racine carrée
va trop profond~; la loi de plaque donnerait un ressaut plus petit.)
\end{solution}

\begin{solution}{exo:b1:rifts-ridges-margins:10}
(a) Avec des aires égales pour chaque âge, la profondeur moyenne est la moyenne de $d(t)$ sur
$[0, A]$~:
\[
  \frac1A\int_0^A\bigl(d_r + C\sqrt t\bigr)\mathrm dt = d_r + \frac CA\cdot\frac23A^{3/2}
  = d_r + \frac23C\sqrt A.
\]
(b) Pour $A = \qty{180}{\Ma}$~: $2500 + \frac23\times350\times13.416 =
\qty{5630}{m}$~; pour \qty{120}{\Ma}~: $2500 + \frac23\times350\times10.954 =
\qty{5056}{m}$, \qty{574}{m} moins profond. (c) Soit $r$ la montée de la mer et $s = r\rho_w/(\rho_m - \rho_w)$
l'enfoncement du plancher chargé. L'eau, de volume inchangé,
doit remplir la même profondeur sur la même aire, donc $r + s = \qty{574}{m}$,
$r\,\rho_m/(\rho_m - \rho_w) = 574$, et $r = 574\times2300/3300 =
\qty{400}{m}$~: des centaines de mètres, par la seule vitesse d'expansion.
\end{solution}

\begin{solution}{exo:b1:rifts-ridges-margins:11}
(a) $\sqrt{x - a} = \sqrt x\sqrt{1 - a/x} \approx \sqrt x\,(1 - a/2x)$ pour
$a \ll x$, donc
$\Delta d \approx C\sqrt{x/u}\cdot a/(2x) = Ca/(2\sqrt{ux})$.
(b) À $x = \qty{450}{km}$~: exact \qty{1205}{m}, approché
$350\times350/(2\sqrt{10.6\times450}) = \qty{887}{m}$~; à \qty{1350}{km}~:
\qty{550}{m} et \qty{512}{m}~; à \qty{3350}{km}~: \qty{334}{m} et
\qty{325}{m}. Le rapport du ressaut approché au ressaut exact ne dépend que de
$r = a/x$~: $(r/2)/(1 - \sqrt{1 - r})$, qui vaut 0.9 pour $r = 0.36$.
L'approximation est bonne à \qty{10}{\percent} près au-delà de $x = 350/0.36 \approx
\qty{970}{km}$, environ trois décalages depuis la dorsale.
\end{solution}

\begin{solution}{exo:b1:rifts-ridges-margins:12}
(a) $d_2 - d_1 = 3671.1 - 3606.8 = \qty{64.3}{m}$ à \qty{10}{\Ma}~;
$4072.8 - 4065.3 = \qty{7.5}{m}$ à 20~; $4415.4 - 4417.0 = \qty{-1.7}{m}$ à
30~; $4707.5 - 4713.6 = \qty{-6.1}{m}$ à 40. (b) $d_2 - d_1$ est positif à 20
et négatif à 30, et les deux lois sont continues~: elles sont égales entre les deux.
Dichotomie~: à 25, $d_2 - d_1 = +0.9$~; à 27.5, $-0.7$~; à 26.25, $+0.01$~; à
26.875, $-0.3$~; à 26.56, $-0.2$~; à 26.41, $-0.08$~; à 26.33, $-0.04$~: la
racine se trouve entre 26.25 et 26.33, à \qty{26.3}{\Ma}, où les deux donnent
\qty{4294}{m}. (c) $d_1' = 175/\sqrt{26.3} = \qty{34.1}{m.\Ma^{-1}}$ et
$d_2' = (3200/62.8)\,e^{-26.3/62.8} = \qty{33.5}{m.\Ma^{-1}}$. Les lois ont
presque la même valeur et la même pente en ce point, donc la courbe raccordée n'a pas de
cassure visible.
\end{solution}

\begin{solution}{pb:b1:rifts-ridges-margins:1}
\textbf{1.} \omterm{def:b1:cross-sections:throw}{Rejet horizontal} $= $ \omterm{def:b1:cross-sections:throw}{rejet vertical}$/\tan$(\omterm{def:b1:geological-maps:attitude}{pendage})~: F1 \qty{346}{m}, F2 \qty{770}{m},
F3 \qty{2014}{m}, F4 \qty{163}{m}, F5 \qty{462}{m}, F6 \qty{1172}{m}. F3, avec
le plus grand \omterm{def:b1:cross-sections:throw}{rejet vertical} et horizontal, est la faille bordière.

\textbf{2.} L'extension est la somme, \qty{4927}{m}, environ \qty{4.9}{km}. Le
rift était large de $60 - 4.9 = \qty{55.1}{km}$ avant la fracturation.

\textbf{3.} La normale à un rift orienté \ang{020} pointe vers \ang{110}~; le
mouvement, vers \ang{105}, fait \ang{5} avec elle. La composante perpendiculaire
est $6.0\cos\ang{5} = \qty{5.98}{mm.\year^{-1}}$.

\textbf{4.} $\qty{4.927}{km}/\qty{5.98}{km.\Ma^{-1}} = \qty{0.82}{\Ma}$.

\textbf{5.} En \qty{3.2}{\Ma} les plaques se sont écartées de
$5.98\times3.2 = \qty{19.1}{km}$ à travers le rift. Les failles en portent
$4.93/19.1 = \qty{26}{\percent}$.

\textbf{6.} Les trois autres quarts, \qty{14.2}{km}, ne sont pas dans les failles de
la coupe. Ils peuvent être repris par des dykes injectés le long de l'axe, qui ouvrent
le rift sans le fracturer, comme en Afar, ou par des failles hors de la coupe.
Avec des dykes larges de \qty{8}{m}, $\qty{14200}{m}/\qty{8}{m} \approx 1780$ dykes sont
nécessaires, un tous les $3.2\times10^6/1780 \approx 1800$ ans en moyenne.

\textbf{7.} À partir de C2An~: $(21.0 + 20.2)/2.581 = 41.2/2.581 =
\qty{15.96}{mm.\year^{-1}}$. À partir de C3n~: $(41.0 + 40.2)/5.235 = 81.2/5.235 =
\qty{15.51}{mm.\year^{-1}}$.

\textbf{8.} $(41.2 - 2.2)/2.581 = \qty{15.11}{mm.\year^{-1}}$ et
$(81.2 - 2.2)/5.235 = \qty{15.09}{mm.\year^{-1}}$~: corrigées, les deux concordent à
\qty{0.02}{mm.an^{-1}} près. La correction compte plus pour le chron le plus jeune, dont la
distance est plus courte.

\textbf{9.} Le déplacement vers l'extérieur décale les deux bords d'un flanc de la même
quantité, donc il s'annule dans leur différence. Sur les deux flancs,
$(81.2 - 41.2)/(5.235 - 2.581) = 40.0/2.654 = \qty{15.07}{mm.\year^{-1}}$.

\textbf{10.} Nord-est~: $(41.0 - 21.0)/2.654 = \qty{7.54}{mm.\year^{-1}}$~;
sud-ouest~: $(40.2 - 20.2)/2.654 = \qty{7.54}{mm.\year^{-1}}$. L'expansion est
symétrique~; la différence de \qty{0.8}{km} entre les flancs, la même pour les deux
chrons, est un décalage constant, comme un axe placé à \qty{0.4}{km} du
centre.

\textbf{11.} $\cos\Delta = \sin\ang{30.9}\sin\ang{18.0} +
\cos\ang{30.9}\cos\ang{18.0}\cos\ang{16.4} = 0.15870 + 0.78287 = 0.94157$, donc
$\Delta = \ang{19.69}$. Avec $\omega R = 0.403\times(\pi/180)\times6371 =
\qty{44.81}{mm.\year^{-1}}$, la vitesse est $44.81\sin\ang{19.69} =
\qty{15.10}{mm.\year^{-1}}$, les vitesses corrigées des questions 8 et 9.

\textbf{12.} $\cos\Delta = \sin\ang{30.9}\sin\ang{26.0} +
\cos\ang{30.9}\cos\ang{26.0}\cos\ang{11.4} = 0.22512 + 0.75600 = 0.98112$, donc
$\Delta = \ang{11.15}$ et la vitesse est $44.81\sin\ang{11.15} =
\qty{8.7}{mm.\year^{-1}}$.

\textbf{13.} Les prédictions, 15.1 et \qty{8.7}{mm.an^{-1}}, correspondent aux
15 et \qty{9}{mm.an^{-1}} mesurés. La vitesse est $\omega R\sin\Delta$, et le
pôle se trouve au nord-ouest de la mer Rouge, dans le nord-est de la Libye~: plus on va vers le sud, plus on est loin
du pôle, plus l'ouverture est rapide.

\textbf{14.} $15.09\times5.235 = \qty{79}{km}$ de plancher océanique.

\textbf{15.} $5000/120.6 = \qty{41}{km.\Ma^{-1}}$, \qty{41}{mm.an^{-1}} en
moyenne.

\textbf{16.} $5000/15.1 = \qty{331}{\Ma}$ à la vitesse du centre~;
$5000/8.7 = \qty{575}{\Ma}$ à la vitesse du nord.

\textbf{17.} En s'ouvrant autour d'un pôle voisin, la mer Rouge s'élargit comme une paire de
ciseaux~: sa largeur croît proportionnellement à $\sin\Delta$, et elle reste un coin
qui se rétrécit vers le pôle. À \qty{8.7}{mm.an^{-1}}, la partie nord a besoin de
$79/8.7 = \qty{9.1}{\Ma}$ pour faire les \qty{79}{km} de plancher océanique que la
partie centrale a maintenant.

\textbf{18.} À \qty{15}{mm.an^{-1}} l'axe de la mer Rouge est une dorsale ultralente~; une
vallée axiale entre des failles normales est attendue.

\textbf{19.} $2500 + 350\times1 = \qty{2850}{m}$;
$2500 + 350\sqrt{10} = \qty{3607}{m}$; $2500 + 350\sqrt{40} = \qty{4714}{m}$.

\textbf{20.} $t = (1500/350)^2 = \qty{18.4}{\Ma}$.

\textbf{21.} Pendant les \qty{10}{\Ma} premiers il s'enfonce de $350\sqrt{10} = \qty{1107}{m}$~;
entre 70 et \qty{80}{\Ma}, $350(\sqrt{80} - \sqrt{70}) = 350\times(8.944 -
8.367) = \qty{202}{m}$, cinq fois moins sur la même durée.

\textbf{22.} $\mathrm dd/\mathrm dt = C/(2\sqrt t) = 175/\sqrt{80} =
\qty{19.6}{m.\Ma^{-1}}$, soit environ \qty{0.02}{mm} par an.

\textbf{23.} $490/\sqrt{80} = \qty{54.8}{mW.m^{-2}}$.

\textbf{24.} Loi en racine carrée~: $2500 + 350\sqrt{80} = 2500 + 350\times8.944 =
\qty{5630}{m}$. Loi de plaque~: $6400 - 3200\,e^{-80/62.8} = 6400 - 3200\times0.2798
= \qty{5505}{m}$. Elles diffèrent de \qty{126}{m}. À \qty{80}{\Ma} les deux lois
s'appliquent encore~: la loi en racine carrée jusqu'à environ \qty{80}{\Ma}, la loi de plaque à partir d'environ
\qty{20}{\Ma}~; au-delà, la loi en racine carrée va trop profond.

\textbf{25.} La loi en racine carrée donne \qty{5630}{m}~: au dixième de
kilomètre près, le plancher du nouvel océan se trouve à une profondeur de
\textbf{\qty[text-series-to-math=true]{5.6}{km}} quand il a 80 millions d'années
(la loi de plaque donne \qty{5.5}{km}).
\end{solution}
